% Начинаем с новой страницы, оформляем как приложение или отдельный документ
\newpage
\begin{flushright}
    {\bfseries ПРИЛОЖЕНИЕ Е} \\
    (обязательное)
\end{flushright}

\begin{center}
    {\bfseries ПЕРЕЧЕНЬ ЭЛЕМЕНТОВ К СХЕМЕ ЭЛЕКТРИЧЕСКОЙ ПРИНЦИПИАЛЬНОЙ} \\
    \vspace{0.2cm}
    БГУИР ДП 1-53 01 07 001 Э3
\end{center}

\addcontentsline{toc}{section}{Приложение Е. Перечень элементов}

\vspace{0.5cm}

\begin{center}
\begin{tabular}{|p{2.5cm}|p{8cm}|p{1.5cm}|p{3cm}|}
\hline
\centering \textbf{Поз. обозначение} & \centering \textbf{Наименование} & \centering \textbf{Кол.} & \textbf{Примечание} \\ \hline
& \underline{Резисторы} & & \\ 
R1, R2 & С2-23-0,125-4,7 кОм $\pm$ 5\% & 2 & \\ \hline
R3 & С2-23-0,125-10 кОм $\pm$ 5\% & 1 & \\ \hline
& \underline{Конденсаторы} & & \\ 
C1 & К10-17-1а-Н90-0,1 мкФ & 1 & \\ \hline
C2 & К50-35-16В-100 мкФ & 1 & \\ \hline
& \underline{Микросхемы} & & \\ 
DD1 & Модуль ESP32-WROOM-32E & 1 & Espressif \\ \hline
DA1 & Стабилизатор AMS1117-3.3 & 1 & \\ \hline
& \underline{Прочие изделия} & & \\ 
U1 & Датчик SHT31-D & 1 & Sensirion \\ \hline
XP1 & Разъем USB-C 16-pin & 1 & \\ \hline
VD1 & Светодиод L-53ID-5V & 1 & Kingbright \\ \hline
\end{tabular}
\end{center}

\vspace{1cm}

\begin{flushleft}
    \begin{tabular}{p{6cm} p{4cm} l}
        Разработал & \underline{\hspace{3cm}} & П. П. Петров \\
        Проверил (руководитель) & \underline{\hspace{3cm}} & С. С. Сидоров \\
        Нормоконтроль & \underline{\hspace{3cm}} & Н. Н. Контрольный
    \end{tabular}
\end{flushleft}

\vfill
\begin{center}
    Минск 2026
\end{center}

\newpage